\chapter{Direção e Board da equipa}

\section{Mandato}
\begin{enumerate}
  \item Os elementos da Board e da Direção estão munidos de mandato enquanto se encontrarem no exercício de funções.
  \item Considera-se o exercício de funções o período de tempo iniciado pela tomada de posse da Board ou pela fixação da Direção, com duração máxima de um ano.
  \item A tomada de posse da Board requer uma eleição legítima e que os eleitos cumpram com o presente Regulamento Geral.
  \item A manutenção do Mandato exige o cumprimento do Regulamento Geral, por parte dos membros em funções.
\end{enumerate}

\section{Eleições}
\begin{enumerate}
  \item A Board da FS FEUP é eleita por voto secreto dos membros ativos da equipa, em âmbito de RGER, de acordo com o calendário eleitoral.
  \item Consideram-se membros ativos, para efeitos eleitorais, os membros que integrem a equipa à data da votação e que não tenham renunciado, sido suspensos ou excluídos nos termos do presente Regulamento.
  \item A votação visa eleger a lista que tomará posse enquanto futura Board da equipa.
  \item A Board eleita assume a responsabilidade de definir e fixar a Direção da equipa, nos termos do presente Regulamento.
  \item O processo eleitoral é organizado pela Board em funções, que assume, para este efeito, a função de Comissão Eleitoral.
  \item O processo eleitoral deve ter início até 15 de maio e terminar até 30 de junho, salvo impossibilidade devidamente fundamentada e comunicada à equipa.
  \item Caso a Board não assuma a função de Comissão Eleitoral ou não dê início ao processo eleitoral até à data estipulada, qualquer membro ativo da equipa, ou conjunto de até cinco membros ativos, pode propor-se para constituir uma Comissão Eleitoral.
  \item A Comissão Eleitoral é constituída na primeira RGER em que seja apresentada uma proposta nesse sentido, assumindo a responsabilidade de definir, organizar e comunicar os prazos e procedimentos do processo eleitoral.
  \item Caso vários membros ou grupos se proponham para constituir a Comissão Eleitoral na mesma RGER, deve ser procurado consenso entre os proponentes, podendo estes fundir propostas ou retirar candidaturas, mantendo sempre o limite máximo de cinco membros.
  \item Não havendo consenso entre propostas concorrentes para Comissão Eleitoral, a composição da Comissão Eleitoral deve ser decidida por votação dos membros presentes em RGER.
  \item Os membros que integrem a Comissão Eleitoral não podem integrar qualquer lista candidata às eleições que lhes caiba organizar.
  \item Quando constituída, a Comissão Eleitoral assume todas as responsabilidades eleitorais que caberiam à Board em situação normal, devendo garantir a neutralidade, transparência e regularidade do processo.
  \item O calendário eleitoral é, de forma sucessiva, composto por:
  \begin{enumerate}
    \item Período de apresentação de listas, com duração mínima de 14 dias.
    \item Período de reflexão, com duração mínima de 7 dias.
    \item Votação.
    \item Período de transição.
    \item Empossamento.
  \end{enumerate}
  \item O processo eleitoral ocorre, anualmente, em datas assinaladas pela Board em funções ou pela Comissão Eleitoral, devendo ser anunciado até uma semana antes do início do calendário eleitoral.
  \item O ato eleitoral deve realizar-se até 30 de junho e, sempre que possível, com pelo menos dois meses de antecedência relativamente ao fim do mandato da Board em funções.
  \item A apresentação formal das listas deve ocorrer em RGER semanal durante o período de apresentação de listas.
  \item Deve realizar-se uma RGER no último dia do período de apresentação de listas ou, se tal não for possível, no dia útil imediatamente anterior.
  \item A votação deve ocorrer em RGER convocada para o efeito.
  \item Caso a RGER necessária para apresentação de listas ou votação não seja convocada pelo Team Leader, a Comissão Eleitoral tem competência para a convocar exclusivamente para esse efeito.
  \item Qualquer membro ativo da equipa tem direito a apresentar uma lista válida.
  \item Uma lista, para ser válida, deve preencher os requisitos:
  \begin{enumerate}
    \item Ser composta, no mínimo, por candidatos aos cargos de Team Leader e Chief Engineer.
    \item Identificar nominalmente os membros candidatos e os cargos da Board a que se propõem.
    \item Ser apresentada de acordo com o calendário eleitoral, à Board em funções ou à Comissão Eleitoral.
    \item Ser apresentada formalmente em RGER, para efeitos de oficialização.
    \item Ter subscrição escrita de um número mínimo correspondente a 10\% dos membros ativos que integram a equipa, externos à lista.
    \item Incluir declaração de aceitação assinada por todos os membros candidatos aos cargos nela indicados.
    \item Caso inclua cargos de Board não previstos no presente Regulamento, apresentar em anexo uma descrição clara das respetivas responsabilidades, âmbito de atuação e enquadramento na estrutura da equipa.
  \end{enumerate}
  \item Um membro pode integrar mais do que uma lista, desde que declare expressamente a aceitação da candidatura em cada uma delas.
  \item Um membro não pode subscrever mais do que uma lista como apoiante.
  \item Não podem existir duas listas exatamente iguais, considerando-se iguais as listas que apresentem os mesmos membros para os mesmos cargos.
  \item Duas listas compostas pelas mesmas pessoas, mas com distribuição diferente de cargos, são consideradas listas distintas.
  \item Cada lista deve escolher uma letra pela qual será identificada durante o processo eleitoral.
  \item A atribuição de letras segue a ordem de oficialização das listas.
  \item Caso todas as letras do alfabeto da língua portuguesa já tenham sido atribuídas, podem ser utilizadas combinações com mais do que uma letra.
  \item Caso uma lista que inclua cargos de Board não previstos no presente Regulamento vença a eleição, esses cargos consideram-se automaticamente integrados no presente Regulamento a partir do empossamento, devendo a descrição anexa à lista constituir a redação aplicável até à sua incorporação formal no texto consolidado.
  \item Consideram-se votos válidos os votos expressos numa lista candidata ou em branco, desde que não contenham elementos que comprometam a sua interpretação ou o segredo de voto.
  \item Considera-se voto branco o voto que não expresse preferência por qualquer lista candidata.
  \item Considera-se voto nulo o voto que não permita identificar de forma inequívoca a intenção do votante, contenha rasuras ou elementos identificativos, ou viole o formato definido para a votação.
  \item A eleição é decidida por maioria simples dos votos válidos expressos em listas candidatas, não sendo os votos brancos considerados para efeitos de apuramento da lista vencedora.
  \item Caso exista apenas uma lista válida, esta apenas será eleita se os votos favoráveis nessa lista forem superiores aos votos brancos.
  \item Durante o momento de contagem dos votos, é permitida a presença de um representante de cada lista candidata, exclusivamente para efeitos de observação.
  \item Os representantes das listas presentes na contagem não podem manusear votos, intervir na contagem ou interferir com o trabalho da Comissão Eleitoral.
  \item Os resultados eleitorais devem ser comunicados à equipa após a contagem dos votos.
  \item Até uma semana após findado o processo de votação e anúncio dos resultados, existe um período de protesto.
  \item Um protesto deve ser fundamentado, conter uma proposta de resolução construtiva e ser apresentado por escrito à Board em funções ou à Comissão Eleitoral.
  \item O protesto pode levar a uma renovação do calendário eleitoral, com possibilidade de nova votação, por decisão da Board em funções ou da Comissão Eleitoral.
  \item O sistema eleitoral funciona a duas voltas:
  \begin{enumerate}
    \item Não havendo numa primeira volta o cumprimento da maioria simples dos votos válidos expressos em listas candidatas para nenhuma lista, é necessária uma segunda volta.
    \item Integrarão numa segunda volta as duas listas com maior número de votos na primeira volta.
    \item O ato de votação da segunda volta deverá ocorrer num prazo de até uma semana da primeira volta.
    \item Em caso de empate na segunda volta, o Team Leader em funções tem direito a voto de qualidade.
  \end{enumerate}
  \item Caso não existam listas válidas, a equipa deve reunir em RGER para definir o seu futuro, podendo aprovar por maioria simples dos membros presentes uma solução transitória, a abertura de novo processo eleitoral extraordinário ou o início de processo de dissolução.
  \item Caso seja aberto novo processo eleitoral extraordinário, este deve respeitar os prazos mínimos de apresentação de listas e reflexão previstos no presente Regulamento.
\end{enumerate}

\section{Definição da Direção pela Board Eleita}
\begin{enumerate}
  \item A Board eleita deve definir e fixar a Direção da equipa no prazo máximo de 25 dias úteis após a eleição.
  \item A Direção definida pela Board eleita deve incluir, obrigatoriamente, pelo menos quatro Department Leaders.
  \item A Direção pode ainda incluir outros membros da equipa nomeados pela Board, mesmo que estes não exerçam funções de liderança departamental.
  \item A Direção apenas se considera validamente definida quando inclua, no mínimo, a Board eleita e quatro Department Leaders.
  \item A Direção final deve ser apresentada em RGER após a sua definição.
  \item Até à apresentação da nova Direção, mantêm-se provisoriamente os Departamentos e responsáveis cessantes, salvo decisão fundamentada da Board eleita.
  \item Compete à Board eleita definir os Departamentos técnicos nos quais a equipa se subdivide e nomear os respetivos Department Leaders.
  \item Para a seleção dos Department Leaders, a Board deve auscultar todos os membros da equipa, de modo a realizar uma decisão informada.
  \item A auscultação deve ser aberta a todos os membros e pode ser realizada por formulário, reunião individual, reunião em grupo ou outro meio definido pela Board.
  \item A auscultação deve procurar recolher, sempre que aplicável, preferências dos membros, disponibilidade, experiência, sugestões de liderança e necessidades técnicas ou organizacionais da equipa.
  \item Em caso de encerramento de um Departamento, a Board deverá procurar alocar os membros desse Departamento a um novo Departamento, tendo em conta as suas preferências e as necessidades da equipa.
  \item Em caso de formação de um novo Departamento, a liderança do mesmo pode ser sugerida no momento de auscultação a todos os membros.
  \item A seleção de membros para um novo Departamento é da responsabilidade da Board, podendo seguir diferentes metodologias, nomeadamente nomeação individual, candidaturas internas, candidaturas externas ou qualquer outro processo de seleção decidido pela Board.
  \item A criação de um novo Departamento não pode prejudicar o direito dos membros a manterem-se no seu Departamento atual, salvo caso de encerramento desse Departamento.
  \item As alterações à estrutura departamental devem ser comunicadas à equipa e devidamente fundamentadas.
  \item Caso a Board eleita não defina uma Direção válida no prazo previsto, é aberto novo processo eleitoral extraordinário.
  \item O processo eleitoral extraordinário deve respeitar os prazos mínimos de apresentação de listas e reflexão previstos no presente Regulamento.
  \item Caso não existam candidaturas no processo eleitoral extraordinário, a equipa deve reunir em RGER para definir o seu futuro, podendo aprovar por maioria simples dos membros presentes uma solução transitória, a abertura de novo processo eleitoral ou o início de processo de dissolução.
\end{enumerate}

\section{Tomada de Posse e Reconfigurações}
\begin{enumerate}
  \item A tomada de posse formal ocorre após a última competição da época, até uma semana após o início do ano letivo.
  \item Em caso de eleição extraordinária, a Board eleita toma posse logo que possível após a validação dos resultados eleitorais, não sendo aplicável o calendário ordinário de tomada de posse.
  \item A Board eleita em processo extraordinário deve definir e apresentar a Direção nos mesmos prazos e termos aplicáveis ao processo eleitoral ordinário.
  \item A Board não deve ser reconfigurada durante o mandato, salvo em caso de necessidade devidamente fundamentada.
  \item A Board não pode ser reconfigurada quando tal afete o Team Leader ou corresponda a alteração superior a 25\% dos seus membros inicialmente eleitos, arredondando por excesso.
  \item Em Boards compostas por apenas dois membros, a substituição de qualquer membro implica novo processo eleitoral.
  \item A renúncia, destituição ou exclusão do Team Leader implica a abertura de novo processo eleitoral.
  \item A Direção pode ser reconfigurada durante o mandato, desde que tal não afete mais de 25\% dos cargos inicialmente fixados, arredondando por excesso.
  \item Qualquer reconfiguração da Direção após a sua fixação inicial carece de aprovação favorável da maioria dos restantes membros da Direção.
  \item Qualquer reconfiguração que não cumpra os requisitos acima mencionados requer a abertura de novo processo eleitoral.
\end{enumerate}
